\section{Formalisme Exact Potential Game}

\subsection{Aksioma Permainan Potensial pada Sistem Finansial}

Teori permainan dalam sistem finansial memodelkan interaksi antar-agen rasional yang berusaha memaksimalkan utilitasnya dalam kondisi ketidakpastian. Untuk mengisi parameter dalam model QUBO dan Hamiltonian Ising yang telah diturunkan pada Subbab 2.3, digunakan pendekatan \textit{Exact Potential Game} (EPG). Dalam konteks optimasi portofolio, setiap aset $i$ dipandang sebagai pemain yang memilih strategi biner $x_i \in \{0, 1\}$. Sistem pemilihan aset ini diklasifikasikan sebagai EPG jika terdapat fungsi potensial global $\Phi(\mathbf{x})$ sedemikian sehingga perubahan utilitas marginal setiap pemain selaras dengan perubahan nilai potensial tersebut \cite{monderer_potential_1996}:
\begin{equation}
\Phi(x_i, \mathbf{x}_{-i}) - \Phi(x'_i, \mathbf{x}_{-i}) = u_i(x_i, \mathbf{x}_{-i}) - u_i(x'_i, \mathbf{x}_{-i})
\label{eq:potential_game_condition}
\end{equation}
Kondisi pada persamaan (\ref{eq:potential_game_condition}) menjamin bahwa pencarian ekuilibrium Nash dalam dinamika pasar dapat direduksi menjadi masalah minimisasi energi pada struktur Hamiltonian Ising.

Bukti bahwa model Markowitz merupakan \textit{Exact Potential Game} dilakukan melalui dekomposisi $\Phi(\mathbf{x})$ untuk memisahkan kontribusi variabel keputusan $x_i$ dari variabel pemain lainnya $\mathbf{x}_{-i}$. Dengan memanfaatkan sifat simetri matriks kovariansi $\mathbf{\Sigma}$ di mana $\sigma_{ij} = \sigma_{ji}$ untuk seluruh pasangan $(i, j)$:
\begin{equation}
\mathbf{\Sigma} = 
\begin{pmatrix}
\sigma_{11} & \sigma_{12} & \cdots & \sigma_{1N} \\
\sigma_{21} & \sigma_{22} & \cdots & \sigma_{2N} \\
\vdots & \vdots & \ddots & \vdots \\
\sigma_{N1} & \sigma_{N2} & \cdots & \sigma_{NN}
\end{pmatrix}
\label{eq:matrix_sigma}
\end{equation}
fungsi potensial diekspansi untuk mengisolasi interaksi lokal pada indeks $i$ sebagai berikut:
\begin{equation}
\Phi(\mathbf{x}) = \mu_i x_i - \frac{\gamma}{2} \left( \sigma_{ii} x_i^2 + 2 \sum_{j \neq i} \sigma_{ij} x_i x_j \right) + C_{-i}
\label{eq:potential_decomposition}
\end{equation}
Suku $C_{-i}$ pada (\ref{eq:potential_decomposition}) memuat seluruh komponen energi yang independen terhadap $x_i$. Mengingat properti peubah biner $x_i^2 = x_i$, faktorisasi variabel keputusan menghasilkan relasi utilitas individual $u_i(x_i, \mathbf{x}_{-i})$ yang eksplisit:
\begin{equation}
\Phi(\mathbf{x}) = x_i \underbrace{\left( \mu_i - \frac{\gamma}{2} \sigma_{ii} - \gamma \sum_{j \neq i} \sigma_{ij} x_j \right)}_{u_i(x_i, \mathbf{x}_{-i})} + C_{-i}
\label{eq:utility_relation}
\end{equation}
Identitas $\Phi(1, \mathbf{x}_{-i}) - \Phi(0, \mathbf{x}_{-i}) = u_i(1, \mathbf{x}_{-i}) - u_i(0, \mathbf{x}_{-i})$ mengonfirmasi eksistensi fungsi potensial eksak, sehingga pencarian ekuilibrium Nash dapat direduksi menjadi pencarian status energi terendah (\textit{ground state}) sistem \cite{monderer_potential_1996}.

\subsection{Konstruksi Utilitas Individual dan Matriks \textit{Payoff}}

Utilitas individual pemain $i$ didefinisikan dengan mengintegrasikan parameter imbal hasil strategis $\tilde{\mu}_i$ dan risiko interaksi yang telah dimodifikasi oleh korelasi non-linear. Formulasi utilitas ini mencakup ekspektasi keuntungan personal, risiko varians individual, serta penalti akibat redundansi informasi dengan aset lain dalam portofolio:
\begin{equation}
u_i(x_i, \mathbf{x}_{-i}) = x_i \left( \tilde{\mu}_i - \frac{\gamma}{2} \sigma_{ii} - \gamma \sum_{j \neq i} \tilde{\sigma}_{ij} x_j \right)
\label{eq:individual_utility_detailed}
\end{equation}
di mana $\gamma$ merupakan koefisien penghindaran risiko yang mengatur bobot relatif antara keuntungan dan volatilitas. Struktur utilitas pada persamaan (\ref{eq:individual_utility_detailed}) memungkinkan setiap agen untuk mengevaluasi kontribusi marginalnya terhadap portofolio kolektif secara otonom namun tetap selaras dengan tujuan sistemik.

Untuk memahami dinamika aljabar dari integrasi parameter strategis tersebut, kita tinjau interaksi antara dua aset yang direpresentasikan sebagai Pemain 1 dan Pemain 2 dalam \textit{normal-form game}. Struktur \textit{payoff} dalam kondisi ini mencerminkan bagaimana korelasi informasional $\tilde{\sigma}_{12}$ secara langsung mereduksi utilitas individual saat kedua aset dipilih secara simultan, sebagaimana disajikan dalam Tabel \ref{tab:payoff_matrix_game}.

\begin{table}[h!]
\centering
\caption{Matriks Payoff $(u_1, u_2)$ untuk Sistem Dua Aset dengan Integrasi Strategis}
\label{tab:payoff_matrix_game}
\begin{tabular}{|c|c|c|}
\hline
\textbf{Pemain 1 \textbackslash Pemain 2} & \textbf{Keluar ($x_2 = 0$)} & \textbf{Masuk ($x_2 = 1$)} \\ \hline
\textbf{Keluar ($x_1 = 0$)} & $(0, 0)$ & $(0, \tilde{\mu}_2 - \frac{\gamma}{2} \sigma_{22})$ \\ \hline
\textbf{Masuk ($x_1 = 1$)} & $(\tilde{\mu}_1 - \frac{\gamma}{2} \sigma_{11}, 0)$ & $(P_{11}, P_{22})$ \\ \hline
\end{tabular}
\end{table}

Pada kondisi di mana kedua pemain memutuskan untuk masuk ke dalam portofolio ($x_1 = 1, x_2 = 1$), \textit{payoff} untuk Pemain 1 secara spesifik didefinisikan sebagai $P_{11} = \tilde{\mu}_1 - \frac{\gamma}{2} \sigma_{11} - \gamma \tilde{\sigma}_{12}$. Hal ini menunjukkan bahwa integrasi redundansi informasi melalui $\tilde{\sigma}_{12}$ memberikan penalti tambahan pada utilitas jika kedua aset memiliki ketergantungan yang tinggi. Penalti strategis ini merupakan mekanisme kunci dalam EPG untuk mendorong diversifikasi aset secara alami melalui mekanisme insentif individual.

\subsection{Ekuivalensi Ekuilibrium Nash dan \textit{Ground State}}

Keunggulan utama dari formulasi \textit{Exact Potential Game} adalah kepastian konvergensi menuju ekuilibrium Nash melalui proses maksimisasi fungsi potensial $\Phi(\mathbf{x})$. Dalam representasi mekanika statistik, ekuilibrium ini bersesuaian tepat dengan status \textit{ground state} dari Hamiltonian Ising yang telah diturunkan. Dengan memetakan utilitas ke dalam lanskap energi, fenomena \textit{herding} atau sinkronisasi keputusan investor dapat dianalisis sebagai transisi fase dalam sistem magnetik sintetis.

Penyelarasan ini memberikan landasan matematis yang kokoh bagi penggunaan algoritma variasional kuantum, karena menjamin bahwa solusi optimal secara global juga merupakan titik kesetimbangan strategis yang stabil bagi seluruh partisipan pasar \cite{menezes_ex_2025}. Hasil akhir dari dekomposisi utilitas dan potensi ini memvalidasi bahwa setiap penurunan energi dalam sistem Ising ekuivalen dengan peningkatan utilitas agen menuju ekuilibrium yang stabil. Oleh karena itu, penggunaan formalisme EPG tidak hanya memberikan interpretasi ekonomi pada parameter fisik Hamiltonian, tetapi juga memastikan integritas logis antara perilaku mikro agen dan optimalitas makro sistem.